% TOP_PROBLEM: {"id":30007625,"problem_number":"LOCAL-30007625","title":"Long-run reshoring effects","category_id":21,"category":{"id":21,"name":"economics","display_name":"Economics","description":"Open economic theory statements, published research agendas, and proposed scoped specifications; question provenance and historical priority are qualified per record.","slug":"economics","order_index":21,"created_at":"2026-10-08T00:00:00Z"},"status":"research_question","external_url":null,"legacy_ids":[],"published":false,"statement":"In the explicitly specified setting: When do tariffs and trade restrictions actually rebuild domestic capacity rather than divert trade or raise downstream costs? The mathematical statement above fixes the target, assumptions and scope of this catalogue formulation.","economics":{"original_id":114,"field":"Trade, globalization and geoeconomics","kind":"empirical","formalization_basis":"adapted_research_question","source_status":"research_agenda","priority_status":"earliest_found","priority_note":"This is the earliest explicit statement verified in this audit; first priority for the full catalogue formulation is not established.","earliest_verified_reference":{"authors":["Bank of England"],"title":"Bank of England Agenda for Research, 2025–2028","year":2026,"url":"https://www.bankofengland.co.uk/research/bank-of-england-agenda-for-research","doi":null,"locator":"Interconnected World / Reconfigured trade, paragraph 1","evidence_summary":"Requests protectionism effects on domestic economy."},"current_reference":{"authors":["Pablo D. Fajgelbaum","Amit Khandelwal"],"title":"Tariffs in 2025: Short-Run Impacts on the U.S. Economy","year":2026,"url":"https://www.brookings.edu/wp-content/uploads/2026/03/1_Fajgelbaum-Khandelwal_unembargoed.pdf","doi":"10.3386/w35064","locator":"Abstract, PDF page 2; Section 5.2.2, final paragraph, PDF page 30; Section 6.4, PDF page 44","evidence_summary":"Finds terms-of-trade evidence inconclusive; longer adjustment and manufacturing outcomes remain uncertain."},"formalization_note":"The source supplies a broad research agenda about protectionism and domestic effects. Durable physical-capacity measurement and reshoring counterfactuals are proposed. The 2026 current paper explicitly describes reshoring and manufacturing outcomes as uncertain; reduced targeted-country imports do not by themselves establish capacity gains."},"statusQualification":"Published question or research agenda verified at the cited locator. The mathematical specification is adapted for this catalogue and includes additional assumptions; source publication does not establish first-ever priority or certify this exact model remains unsolved. The source supplies a broad research agenda about protectionism and domestic effects. Durable physical-capacity measurement and reshoring counterfactuals are proposed. The 2026 current paper explicitly describes reshoring and manufacturing outcomes as uncertain; reduced targeted-country imports do not by themselves establish capacity gains.","statusReviewedAt":"2026-10-08"}
% FIRST_POSTED: 2026-10-08
% ENTRY_KIND: statement_only
% !TeX program = lualatex
\documentclass[11pt]{article}
\usepackage{fontspec}
\usepackage[a4paper,margin=25mm]{geometry}
\usepackage{amsmath,amssymb,mathtools}
\usepackage{xurl}
\usepackage[hidelinks]{hyperref}
\newcommand{\sref}[2]{\hyperlink{source:#1:#2}{[S#2]}}
\setlength{\emergencystretch}{3em}
\begin{document}
% problemId:30007625
\section*{Long-run reshoring effects}\label{30007625}
\subsection{Definitions and mathematical statement}
Fix industries \(j\), countries and a tariff or sourcing-restriction policy path \(a\). Let \(K^d_{j,t+h}(a)\) be productive domestic capacity, defined by a prespecified engineering or production metric rather than nominal sales. Let \(V^d_{j,t+h}(a)\) be real domestic value added and \(C^{down}_{t+h}(a)\) downstream production costs. Define reshoring effects relative to an otherwise specified baseline path \(0\):
\[\Delta K_j(h)=E[K^d_{j,t+h}(do(a))-K^d_{j,t+h}(do(0))].\]
The task is to identify when positive capacity and value-added changes persist after the initial policy shock, alongside employment, productivity and downstream-cost effects. Specify the long horizon \(h\), policy permanence and capacity utilization. Distinguish relocation to the home country from import diversion through third countries, transshipment or increased imported inputs embedded in domestic assembly. Account for retaliation, foreign investment, exchange rates and endogenous entry. Identify or bound counterfactual capacity when policy selects industries already expected to expand or decline. Compare local adjustment estimates with national general-equilibrium outcomes. Evidence of reduced direct imports from a targeted partner does not establish durable domestic productive-capacity growth or a positive net welfare effect.

\par\textbf{Formulation and scope.} The source supplies a broad research agenda about protectionism and domestic effects. Durable physical-capacity measurement and reshoring counterfactuals are proposed. The 2026 current paper explicitly describes reshoring and manufacturing outcomes as uncertain; reduced targeted-country imports do not by themselves establish capacity gains.

\par\textbf{Open-question provenance.} An explicit question or research agenda was verified at the source locator. This does not by itself certify that every later version remains unresolved \sref{30007625}{1}.

\par\textbf{Historical priority.} This is the earliest explicit statement verified in this audit; first priority for the full catalogue formulation is not established.

\subsection{Short English statement}
In the explicitly specified setting: When do tariffs and trade restrictions actually rebuild domestic capacity rather than divert trade or raise downstream costs? The mathematical statement above fixes the target, assumptions and scope of this catalogue formulation.

\subsection{Sources}
\begin{itemize}\raggedright
\item[\textnormal{[S1]}]\hypertarget{source:30007625:1}{} Bank of England. 2026. Bank of England Agenda for Research, 2025–2028. Interconnected World / Reconfigured trade, paragraph 1.
\url{https://www.bankofengland.co.uk/research/bank-of-england-agenda-for-research}.
{\small\textit{Earliest verified question/agenda reference; historical priority qualified below. Requests protectionism effects on domestic economy.}}

\item[\textnormal{[S2]}]\hypertarget{source:30007625:2}{} Pablo D. Fajgelbaum, Amit Khandelwal. 2026. Tariffs in 2025: Short-Run Impacts on the U.S. Economy. Abstract, PDF page 2; Section 5.2.2, final paragraph, PDF page 30; Section 6.4, PDF page 44.
\url{https://www.brookings.edu/wp-content/uploads/2026/03/1_Fajgelbaum-Khandelwal_unembargoed.pdf}.
DOI: 10.3386/w35064.
{\small\textit{Supporting or current literature; see evidence qualification. Finds terms-of-trade evidence inconclusive; longer adjustment and manufacturing outcomes remain uncertain.}}
\end{itemize}
\end{document}
